% Môn: Vật lý
% Lớp: 12
% Chương: Chương I. Vật lí nhiệt
% Bài: L12C1 Bài 4. Nhiệt dung riêng · Tính nhiệt lượng, khối lượng và nhiệt hóa hơi riêng bằng Q = Lm
% ===== Câu 51 =====
% ID: L12C1B4-15-TLN
% Mức: NB
\dangbt{Tính nhiệt lượng, khối lượng và nhiệt hóa hơi riêng bằng Q = Lm}
\begin{ex}
% ID: L12C1B4-15-TLN
% Mức: NB
Cồn bay hơi $20\,\text{g}, L=0,9.10^6\,\text{J}$ /kg. Nhiệt lượng thu theo kJ là bao nhiêu? Chỉ ghi số.
\shortans{18}
\loigiai{
$m=0,020 kg; Q=0,9.10^6.0,020=18 kJ.$
}
\end{ex}

% ===== Câu 55 =====
% ID: L12C1B4-16-TLN
% Mức: TH
\begin{ex}
% ID: L12C1B4-16-TLN
% Mức: TH
\shortans{18}
\loigiai{
Nhiệt lượng mà thiết bị lấy khỏi phòng là $Q = 900 \text{ kJ} = 900 \times 10^3 \text{ J}$.
Khối lượng môi chất bay hơi là $m = 0,50 \text{ kg}$.
Nhiệt hóa hơi riêng $L$ được tính bằng công thức $L = \frac{Q}{m}$.
Thay số vào, ta có:
$L = \frac{900 \times 10^3 \text{ J}}{0,50 \text{ kg}} = 1800 \times 10^3 \text{ J/kg} = 1,8 \times 10^6 \text{ J/kg}$.
Để biểu diễn $L$ theo đơn vị $10^5 \text{ J/kg}$, ta chuyển đổi:
$L = 1,8 \times 10^6 \text{ J/kg} = 18 \times 10^5 \text{ J/kg}$.
Vậy, giá trị của $L$ theo đơn vị $10^5 \text{ J/kg}$ là $18$.
}
\end{ex}

% ===== Câu 59 =====
% ID: L12C1B4-17-TLN
% Mức: TH
\begin{ex}
% ID: L12C1B4-17-TLN
% Mức: TH
\shortans{115}
\loigiai{
Khối lượng nước bay hơi là $m = 50 \text{ g} = 0,050 \text{ kg}$.
Nhiệt hóa hơi riêng của nước là $L = 2,3 \cdot 10^6 \text{ J/kg}$.
Nhiệt lượng mà nước đã lấy từ môi trường để bay hơi được tính bằng công thức:
$Q = m \cdot LQ = 0,050 \text{ kg} \cdot 2,3 \cdot 10^6 \text{ J/kg}Q = 115000 \text{ J}$
Để đổi sang kilojoule (kJ), ta chia cho 1000:
$Q = \frac{115000}{1000} \text{ kJ}Q = 115 \text{ kJ}$
Vậy, nhiệt lượng mà nước đã lấy từ môi trường là $115 \text{ kJ}$.
}
\end{ex}

% ===== Câu 63 =====
% ID: L12C1B4-18-TLN
% Mức: VD
\begin{ex}
% ID: L12C1B4-18-TLN
% Mức: VD
\shortans{30}
\loigiai{
Nhiệt lượng mà dàn lạnh thu được mỗi phút được tính bằng công thức $Q = L \cdot m$.
Trong đó:
$L = 1,5 \cdot 10^5 \text{ J/kg}$ là nhiệt hóa hơi riêng của môi chất.
$m = 0,20 \text{ kg}$ là khối lượng môi chất bay hơi mỗi phút.
Thay số vào công thức, ta có:
$Q = (1,5 \cdot 10^5 \text{ J/kg}) \cdot (0,20 \text{ kg}) = 30000 \text{ J}$.
Để đổi từ Joule (J) sang kilojoule (kJ), ta chia cho 1000:
$Q = \frac{30000}{1000} \text{ kJ} = 30 \text{ kJ}$.
Vậy, nhiệt lượng mà dàn lạnh thu được mỗi phút là $30 \text{ kJ}$.
}
\end{ex}

% ===== Câu 67 =====
% ID: L12C1B4-53-TLN
% Mức: VDC
\begin{ex}
% ID: L12C1B4-53-TLN
% Mức: VDC
\shortans{230}
\loigiai{
Trong quá trình ngưng tụ, hơi nước tỏa ra nhiệt lượng. Nhiệt lượng tỏa ra được tính bằng công thức:
$Q = L \cdot m$
Trong đó:
$L$ là nhiệt hóa hơi riêng của chất lỏng ($L = 2,3 \cdot 10^6 \text{ J/kg}$).
$m$ là khối lượng của hơi nước ($m = 0,10 \text{ kg}$).
Thay số vào công thức, ta có:
$Q = (2,3 \cdot 10^6 \text{ J/kg}) \cdot (0,10 \text{ kg})Q = 2,3 \cdot 10^5 \text{ J}$
Để chuyển đổi nhiệt lượng từ joule (J) sang kilojoule (kJ), ta chia cho $1000$:
$Q = \frac{2,3 \cdot 10^5}{1000} \text{ kJ}Q = 230 \text{ kJ}$
Vậy, nhiệt lượng tỏa ra trong quá trình ngưng tụ là $230 \text{ kJ}$.
}
\end{ex}

% ===== Câu 71 =====
% ID: L12C1B4-54-TLN
% Mức: VDC
\begin{ex}
% ID: L12C1B4-54-TLN
% Mức: VDC
\shortans{200}
\loigiai{
Nhiệt lượng tỏa ra khi hơi nước ngưng tụ là $Q = 460 \text{ kJ} = 460000 \text{ J}$.
Nhiệt hóa hơi riêng của nước là $L = 2,3 \cdot 10^6 \text{ J/kg}$.
Công thức tính nhiệt lượng tỏa ra khi ngưng tụ là $Q = m \cdot L$.
Từ đó, khối lượng hơi nước đã ngưng tụ là $m = \frac{Q}{L}$.
Thay số vào, ta có $m = \frac{460000}{2,3 \cdot 10^6} = \frac{460000}{2300000} = 0,2 \text{ kg}$.
Để đổi sang gam, ta nhân với 1000: $m = 0,2 \cdot 1000 = 200 \text{ g}$.
}
\end{ex}

% ===== Câu 101 =====
% ID: L12C1B4-55-DS
% Mức: NB
\begin{ex}
% ID: L12C1B4-55-DS
% Mức: NB
Để xác định $L$ bằng phép đo $Q$ và m.
\choiceTF
{\True Có thể dùng L=Q/m.}
{\True $Q$ phải là phần năng lượng hữu ích cho hóa hơi.}
{\True m cần đổi về kg nếu tính $L$ theo J/kg.}
{Sai số của $Q$ và m không ảnh hưởng kết quả.}
\loigiai{
\begin{itemchoice}
\item (Đúng) Có thể dùng L=Q/m.
\item (Đúng) $Q$ phải là phần năng lượng hữu ích cho hóa hơi. (Vì): a phát biểu đầu là điều kiện tính; đại lượng đo có sai số nên $L$ cũng có sai số
\item (Đúng) m cần đổi về kg nếu tính $L$ theo J/kg.
\item (Sai) Sai số của $Q$ và m không ảnh hưởng kết quả.
\end{itemchoice}
}
\end{ex}

% ===== Câu 105 =====
% ID: L12C1B4-56-DS
% Mức: TH
\begin{ex}
% ID: L12C1B4-56-DS
% Mức: TH
Hai mẫu cùng chất ở nhiệt độ sôi có m_A=2m_B.
\choiceTF
{\True $Q_A=2Q_B.$}
{\True $Q_A/Q_B=2.$}
{$L_A=2L_B.$}
{\True Nếu ngưng tụ thì độ lớn nhiệt lượng của $A$ cũng gấp đôi B.}
\loigiai{
\begin{itemchoice}
\item (Đúng) $Q_A=2Q_B.$.
\item (Đúng) $Q_A/Q_B=2.$.
\item (Sai) $L_A=2L_B.$. (Vì): ùng chất nên L_A=L_ $B$; Q=Lm
\item (Đúng) Nếu ngưng tụ thì độ lớn nhiệt lượng của $A$ cũng gấp đôi B.
\end{itemchoice}
}
\end{ex}

% ===== Câu 109 =====
% ID: L12C1B4-57-DS
% Mức: VD
\begin{ex}
% ID: L12C1B4-57-DS
% Mức: VD
Hóa hơi $0,50\,\text{kg}$ nước ở 100 ° $C$, lấy L=2,3.10^ $6\,\text{J}$ /kg.
\choiceTF
{\True $Q=1,15.10^6 J.$}
{\True $Q=1150 kJ.$}
{\True Nếu chỉ hóa hơi $0,25\,\text{kg}$ thì $Q$ giảm một nửa.}
{$Q$ không phụ thuộc khối lượng.}
\loigiai{
\begin{itemchoice}
\item (Đúng) $Q=1,15.10^6 J.$.
\item (Đúng) $Q=1150 kJ.$.
\item (Đúng) Nếu chỉ hóa hơi $0,25\,\text{kg}$ thì $Q$ giảm một nửa.
\item (Sai) $Q$ không phụ thuộc khối lượng.
\end{itemchoice}
}
\end{ex}

% ===== Câu 113 =====
% ID: L12C1B4-58-DS
% Mức: VD
\begin{ex}
% ID: L12C1B4-58-DS
% Mức: VD
Kiểm tra đơn vị trong Q=Lm.
\choiceTF
{\True J/kg nhân kg cho J.}
{\True kJ có thể đổi sang $J$ bằng cách nhân 1000.}
{Gam có thể thay trực tiếp cho kilôgam mà không đổi.}
{\True Nếu $Q$ tính kJ và m kg thì $L$ tính kJ/kg.}
\loigiai{
\begin{itemchoice}
\item (Đúng) J/kg nhân kg cho J.
\item (Đúng) kJ có thể đổi sang $J$ bằng cách nhân 1000.
\item (Sai) Gam có thể thay trực tiếp cho kilôgam mà không đổi.
\item (Đúng) Nếu $Q$ tính kJ và m kg thì $L$ tính kJ/kg.
\end{itemchoice}
}
\end{ex}

% ===== Câu 117 =====
% ID: L12C1B4-59-DS
% Mức: VDC
\begin{ex}
% ID: L12C1B4-59-DS
% Mức: VDC
Một chất có L=1,2.10^ $6\,\text{J}$ /kg; cung cấp 360 kJ chỉ cho hóa hơi.
\choiceTF
{\True Khối lượng hóa hơi là $0,30\,\text{kg}$.}
{\True Khối lượng hóa hơi là $300\,\text{g}$.}
{Nếu hiệu suất chỉ 50% thì cùng năng lượng nguồn vẫn hóa hơi $0,30\,\text{kg}$.}
{\True $L=Q/m.$}
\loigiai{
\begin{itemchoice}
\item (Đúng) Khối lượng hóa hơi là $0,30\,\text{kg}$.
\item (Đúng) Khối lượng hóa hơi là $300\,\text{g}$.
\item (Sai) Nếu hiệu suất chỉ 50% thì cùng năng lượng nguồn vẫn hóa hơi $0,30\,\text{kg}$.
\item (Đúng) $L=Q/m.$.
\end{itemchoice}
}
\end{ex}

% ===== Câu 121 =====
% ID: L12C1B4-60-TL
% Mức: NB
\begin{bt}
% ID: L12C1B4-60-TL
% Mức: NB
Hai chất $A$ và $B$ cùng khối lượng $0,30\,\text{kg}$ có L_A=2,0.10^ $6\,\text{J}$ /kg, L_B=0,9.10^ $6\,\text{J}$ /kg. Tính và so sánh nhiệt lượng hóa hơi.
\loigiai{
Q_A=2,0.10^6.0,30=600 kJ; Q_B=0,9.10^6.0,30=270 kJ. Chất $A$ cần nhiều hơn 330 kJ, tỉ số Q_A/Q_B=2,22.
}
\end{bt}

% ===== Câu 122 =====
% ID: L12C1B4-61-TLN
% Mức: NB
\begin{ex}
% ID: L12C1B4-61-TLN
% Mức: NB
Cung cấp 540 kJ để hóa hơi chất có L=1,8.10^ $6\,\text{J}$ /kg. Tính khối lượng theo gam, chỉ ghi số.
\shortans{300}
\loigiai{
$m=540000/(1,8.10^6)=0,30 kg=300 g.$
}
\end{ex}

% ===== Câu 124 =====
% ID: L12C1B4-62-TL
% Mức: TH
\begin{bt}
% ID: L12C1B4-62-TL
% Mức: TH
Tính nhiệt lượng cần để hóa hơi hoàn toàn $0,75\,\text{kg}$ nước ở 100 ° $C$, lấy L=2,3.10^ $6\,\text{J}$ /kg.
\loigiai{
Áp dụng Q=Lm=2,3.10^6.0,75=1,725.10^ $6\,\text{J}$ =1725 kJ.
}
\end{bt}

% ===== Câu 125 =====
% ID: L12C1B4-74-TLN
% Mức: NB
\begin{ex}
% ID: L12C1B4-74-TLN
% Mức: NB
Q_A=720 kJ dùng hóa hơi $0,40\,\text{kg}$ chất A. Tính L_ $A$ theo đơn vị 10^ $6\,\text{J}$ /kg, chỉ ghi số.
\shortans{01/08/2026}
\loigiai{
$L=720000/0,40=1,8.10^6 J/kg.$
}
\end{ex}

% ===== Câu 127 =====
% ID: L12C1B4-75-TL
% Mức: TH
\begin{bt}
% ID: L12C1B4-75-TL
% Mức: TH
Một mẫu hơi $0,50\,\text{kg}$ ngưng tụ hoàn toàn, L=1,6.10^ $6\,\text{J}$ /kg. Tính nhiệt lượng tỏa ra và nêu dấu của $Q$ nếu quy ước hệ nhận nhiệt dương.
\loigiai{
Độ lớn nhiệt lượng tỏa ra là Lm=1,6.10^6.0,50=8,0.10^ $5\,\text{J}$. Theo quy ước hệ nhận nhiệt dương, hơi tỏa nhiệt nên Q=-8,0.10^ $5\,\text{J}$.
}
\end{bt}

% ===== Câu 128 =====
% ID: L12C1B4-78-TLN
% Mức: TH
\begin{ex}
% ID: L12C1B4-78-TLN
% Mức: TH
Hóa hơi $250\,\text{g}$ chất lỏng cần 200 kJ. Tính $L$ theo đơn vị 10^ $5\,\text{J}$ /kg, chỉ ghi số.
\shortans{8}
\loigiai{
$m=0,25 kg; L=200000/0,25=8,0.10^5 J/kg.$
}
\end{ex}

% ===== Câu 130 =====
% ID: L12C1B4-79-TL
% Mức: VD
\begin{bt}
% ID: L12C1B4-79-TL
% Mức: VD
Một bình mất $0,24\,\text{kg}$ chất lỏng do sôi; năng lượng hữu ích đã truyền là 432 kJ. Xác định nhiệt hóa hơi riêng.
\loigiai{
$L=Q/m=432000/0,24=1,8.10^6 J/kg.$
}
\end{bt}

% ===== Câu 131 =====
% ID: L12C1B4-80-TLN
% Mức: VD
\begin{ex}
% ID: L12C1B4-80-TLN
% Mức: VD
Ngưng tụ $0,12\,\text{kg}$ hơi nước, lấy L=2,3.10^ $6\,\text{J}$ /kg. Tính nhiệt lượng tỏa ra theo kJ, chỉ ghi số.
\shortans{276}
\loigiai{
$Q=Lm=276000 J=276 kJ.$
}
\end{ex}

% ===== Câu 133 =====
% ID: L12C1B4-81-TL
% Mức: VD
\begin{bt}
% ID: L12C1B4-81-TL
% Mức: VD
Chứng minh cách xác định $L$ từ đồ thị $Q$ theo m của cùng một chất ở nhiệt độ sôi.
\loigiai{
Từ Q=Lm, đồ thị $Q$ theo m là đường thẳng qua gốc. Hệ số góc ΔQ/Δm chính là L. Đo nhiều cặp (m,Q), vẽ đường phù hợp và lấy độ dốc giúp giảm ảnh hưởng sai số ngẫu nhiên.
}
\end{bt}

% ===== Câu 134 =====
% ID: L12C1B4-82-TLN
% Mức: VDC
\begin{ex}
% ID: L12C1B4-82-TLN
% Mức: VDC
Hóa hơi $0,18\,\text{kg}$ nước ở nhiệt độ sôi với L=2,3.10^ $6\,\text{J}$ /kg. Tính $Q$ theo kJ, chỉ ghi số.
\shortans{414}
\loigiai{
$Q=2,3.10^6.0,18=414000 J=414 kJ.$
}
\end{ex}

% ===== Câu 136 =====
% ID: L12C1B4-83-TL
% Mức: VDC
\begin{bt}
% ID: L12C1B4-83-TL
% Mức: VDC
Cung cấp 920 kJ để hóa hơi nước ở 100 °C. Lấy L=2,3.10^ $6\,\text{J}$ /kg. Tính khối lượng nước hóa hơi.
\loigiai{
$m=Q/L=920000/(2,3.10^6)=0,40 kg.$
}
\end{bt}

% ===== Câu 137 =====
% ID: L12C1B4-84-TLN
% Mức: VDC
\begin{ex}
% ID: L12C1B4-84-TLN
% Mức: VDC
Một nguồn cung cấp 1,2 MJ hữu ích cho chất có L=2,0.10^ $6\,\text{J}$ /kg. Khối lượng hóa hơi theo gam là bao nhiêu? Chỉ ghi số.
\shortans{600}
\loigiai{
$m=1,2.10^6/(2,0.10^6)=0,60 kg=600 g.$
}
\end{ex}

% ===== Câu 154 =====
% ID: L12C1B4-85-TN
% Mức: NB
\begin{ex}
% ID: L12C1B4-85-TN
% Mức: NB
Hóa hơi hoàn toàn $0,2\,\text{kg}$ chất lỏng ở nhiệt độ sôi, biết L=2,3.10^ $6\,\text{J}$ /kg. Nhiệt lượng cần cung cấp là
\choice
{46 kJ}
{\True 460 kJ}
{4600 kJ}
{460 $J$}
\loigiai{
$Q=Lm=2,3.10^6.0,2=460 kJ.$
}
\end{ex}

% ===== Câu 155 =====
% ID: L12C1B4-86-TN
% Mức: NB
\begin{ex}
% ID: L12C1B4-86-TN
% Mức: NB
Hóa hơi hoàn toàn $0,25\,\text{kg}$ chất lỏng ở nhiệt độ sôi, biết L=2,4.10^ $6\,\text{J}$ /kg. Nhiệt lượng cần cung cấp là
\choice
{60 kJ}
{\True 600 kJ}
{6000 kJ}
{600 $J$}
\loigiai{
$Q=Lm=2,4.10^6.0,25=600 kJ.$
}
\end{ex}

% ===== Câu 156 =====
% ID: L12C1B4-87-TN
% Mức: NB
\begin{ex}
% ID: L12C1B4-87-TN
% Mức: NB
Hóa hơi $0,30\,\text{kg}$ chất lỏng cần 270 kJ. Nhiệt hóa hơi riêng là
\choice
{$9,0.10^4 J/kg$}
{\True $9,0.10^5 J/kg$}
{$9,0.10^6 J/kg$}
{$8,1.10^4 J/kg$}
\loigiai{
$L=Q/m=270000/0,30=9,0.10^5 J/kg.$
}
\end{ex}

% ===== Câu 157 =====
% ID: L12C1B4-169-TN
% Mức: TH
\begin{ex}
% ID: L12C1B4-169-TN
% Mức: TH
Hóa hơi hoàn toàn $0,35\,\text{kg}$ chất lỏng ở nhiệt độ sôi, biết L=2.10^ $6\,\text{J}$ /kg. Nhiệt lượng cần cung cấp là
\choice
{70 kJ}
{\True 700 kJ}
{7000 kJ}
{700 $J$}
\loigiai{
$Q=Lm=2.10^6.0,35=700 kJ.$
}
\end{ex}

% ===== Câu 158 =====
% ID: L12C1B4-170-TN
% Mức: TH
\begin{ex}
% ID: L12C1B4-170-TN
% Mức: TH
Cần 460 kJ để hóa hơi chất lỏng có L=2,3.10^ $6\,\text{J}$ /kg. Khối lượng chất lỏng là
\choice
{$0,02\,\text{kg}$}
{\True $0,2\,\text{kg}$}
{$2\,\text{kg}$}
{$200\,\text{kg}$}
\loigiai{
$m=Q/L=460000/(2,3.10^6)=0,2 kg.$
}
\end{ex}

% ===== Câu 159 =====
% ID: L12C1B4-171-TN
% Mức: TH
\begin{ex}
% ID: L12C1B4-171-TN
% Mức: TH
Hai chất $A$, $B$ có cùng khối lượng hóa hơi; L_A=2L_B. Tỉ số Q_A/Q_ $B$ là
\choice
{01/02/2026}
{1}
{\True 2}
{4}
\loigiai{
Q=Lm nên với cùng m, Q_A/Q_B=L_A/L_B=2.
}
\end{ex}

% ===== Câu 160 =====
% ID: L12C1B4-172-TN
% Mức: VD
\begin{ex}
% ID: L12C1B4-172-TN
% Mức: VD
Hóa hơi hoàn toàn $0,4\,\text{kg}$ chất lỏng ở nhiệt độ sôi, biết L=0,9.10^ $6\,\text{J}$ /kg. Nhiệt lượng cần cung cấp là
\choice
{36 kJ}
{\True 360 kJ}
{3600 kJ}
{360 $J$}
\loigiai{
$Q=Lm=0,9.10^6.0,4=360 kJ.$
}
\end{ex}

% ===== Câu 161 =====
% ID: L12C1B4-173-TN
% Mức: VD
\begin{ex}
% ID: L12C1B4-173-TN
% Mức: VD
Cần 600 kJ để hóa hơi chất lỏng có L=2.10^ $6\,\text{J}$ /kg. Khối lượng chất lỏng là
\choice
{$0,03\,\text{kg}$}
{\True $0,3\,\text{kg}$}
{$3\,\text{kg}$}
{$300\,\text{kg}$}
\loigiai{
$m=Q/L=600000/(2.10^6)=0,3 kg.$
}
\end{ex}

% ===== Câu 162 =====
% ID: L12C1B4-174-TN
% Mức: VDC
\begin{ex}
% ID: L12C1B4-174-TN
% Mức: VDC
Hóa hơi hoàn toàn $0,6\,\text{kg}$ chất lỏng ở nhiệt độ sôi, biết L=1,5.10^ $6\,\text{J}$ /kg. Nhiệt lượng cần cung cấp là
\choice
{90 kJ}
{\True 900 kJ}
{9000 kJ}
{900 $J$}
\loigiai{
$Q=Lm=1,5.10^6.0,6=900 kJ.$
}
\end{ex}

% ===== Câu 163 =====
% ID: L12C1B4-175-TN
% Mức: VDC
\begin{ex}
% ID: L12C1B4-175-TN
% Mức: VDC
Cần 450 kJ để hóa hơi chất lỏng có L=1,5.10^ $6\,\text{J}$ /kg. Khối lượng chất lỏng là
\choice
{$0,03\,\text{kg}$}
{\True $0,3\,\text{kg}$}
{$3\,\text{kg}$}
{$300\,\text{kg}$}
\loigiai{
$m=Q/L=450000/(1,5.10^6)=0,3 kg.$
}
\end{ex}
